\section{Experimental Setup}
\label{sec:setup}

\begin{table}[t]
\centering
\small
\caption{The eight benchmark settings with 50 or more tasks (agents / species / tasks) and the budget $B_1$, the published computation time of RL(s.10) \cite{dai2025heterogeneous}. $N$: range of the number of rollouts that the released policy draws within $B_1$ on our 8 cores.}
\label{tab:settings}
\begin{tabular}{lrrc}
\toprule
Setting & $B_1$ (s) & RL(s.$N$): $N$ & CTAS-D run \\
\midrule
SA-BT 25/5/50 & 13.45 & 154--192 & yes \\
SA-BT 50/5/50 & 19.14 & 144--162 & yes \\
SA-AT 50/5/50 & 29.80 & 140--166 & yes \\
MA-AT 25/5/50 & 16.34 & 157--241 & yes \\
MA-AT 50/5/50 & 23.43 & 135--194 & yes \\
MA-AT 50/5/200 & 75.20 & 40--67 & yes \\
MA-AT 150/10/500 & 560.20 & 32--36 & no \\
MA-AT 150/5/500 & 560.50 & 32--40 & no \\
\bottomrule
\end{tabular}
\end{table}


\paragraph{Instances.} We study the eight settings of the benchmark with 50 or more tasks (Table~\ref{tab:settings}).
As pre-registered, the confirmatory test leaves out the eight 20-task settings, which exact methods reach (CP-SAT
proves every val instance of SA-BT 25/5/20 optimal); the supplement reports three of them descriptively. Instances are generated with the benchmark's generator from disjoint seed ranges into three splits
per setting: \emph{dev} (20 instances), on which ALNS and every competitor were tuned; \emph{val} (20 instances),
used only for dry runs, anytime curves and reference solutions; and \emph{test} (50 instances), run once. On our
test instances the released RL policy reproduces the published mean makespans of RL(g.) and RL(s.10) within half a
standard deviation on all 16 settings of the benchmark, which confirms both the instances and our use of the policy.

\paragraph{Competitors.} Every competitor is compared separately and runs on the same 8 cores as ALNS.
\begin{itemize}
\item \textbf{RL(s.$N$)}: the released policy of \citet{dai2025heterogeneous} in 8 single-threaded processes, each
  sampling rollouts in lockstep batches until the budget is used; the best rollout counts. It draws \RLNlo{} to
  \RLNhi{} times as many rollouts as the published RL(s.10) (Table~\ref{tab:settings}).
\item \textbf{CP-SAT LNS}: LNS with the CP-SAT solver \cite{perron2026cpsat}. Each step frees a window of tasks
  (critical chain, spatial, temporal or random), fixes all other coalitions and re-solves a CP model with one circuit
  constraint per affected agent, with 8 threads per sub-solve, starting from our construction heuristic.
\item \textbf{Parallel CP-SAT LNS}: the same with 8 single-threaded workers that improve one shared incumbent,
  starting from the best of 8 construction streams, so that all cores stay busy.
\item \textbf{CP-SAT model}: the full CP-SAT model of the problem (a circuit per agent over its candidate tasks,
  coverage constraints, symmetry breaking between identical agents) solved by CP-SAT's parallel portfolio of
  complete and LNS workers \cite{perron2023cpsatlp} with 8 threads, hinted with the best of 8 construction streams.
\item \textbf{CTAS-D}: the MILP of \citet{fu2023robust} that the benchmark paper used, which we port from the
  authors' public code to CP-SAT as the MIP solver (no Gurobi licence), with 8 threads and no warm start. The shipped
  Gurobi solutions of the benchmark's 50 released test instances satisfy our model and reproduce their objective
  values, and on ten of them our port finds, in 60\,s, plans 4.1\% shorter than the shipped 600-s Gurobi runs. On the
  500-task settings it found no solution within $B_1$ on dev and used 55--87\,GB of memory, so it is counted as
  failing there and not run.
\item \textbf{Restarts}: 8 independent streams of our randomized constructions (Section~\ref{sec:search}); the best
  plan finished within the budget counts.
\end{itemize}
The sub-solve length of parallel CP-SAT LNS (0.5--5\,s) and the arc restriction of the CP-SAT model (all arcs or the
5--20 nearest tasks) were chosen per setting on dev, where CP-SAT LNS keeps its setting from earlier tuning (2\,s
sub-solves, 10\,s on MA-AT 150/10/500). Every plan is scored by replaying it in the benchmark's simulator; RL is scored
by its own rollout. A run without a plan, or whose plan does not finish all tasks before the simulator's limit,
counts as makespan 200.

\paragraph{Compute.} The budget $B_1$ of a setting is the computation time that the benchmark paper reports for
RL(s.10), and $2B_1$ is its double. All runs used one host with two Intel Xeon Platinum 8480C processors in a container
with a quota of 57.6 CPUs, running at most six lanes of 8 CPUs, one per physical core, with the jobs of all methods
shuffled into the same lanes. Budgets start once the instance is loaded in a warm process (compiled kernels, loaded
network weights) and include construction and model building. A lane starts a job only when at most 40\% of the
container's CPU periods were throttled in the last second; a job that saw more throttling is re-run. A run that
overruns its budget by more than 10\% + 0.25\,s is scored by the best plan it had found by the budget, or as a failure
if it records none. We report the CPU time that each method actually used.

\paragraph{Hypotheses and analysis.} The pre-registered hypotheses are:
\textbf{C1}: on each setting, ALNS on 8 cores at $B_1$ has a lower mean makespan than each competitor on 8 cores at
$B_1$ (46 tests: 8 settings $\times$ 6 competitors, without CTAS-D on the 500-task settings);
\textbf{C1 at $2B_1$}: the same against parallel CP-SAT LNS, the CP-SAT model and restarts at $2B_1$ on the six settings
up to 200 tasks (17 tests); and \textbf{C2}: ALNS on one core for 2\,s is better than each competitor on 8 cores at
$B_1$ (46 tests). For each instance we compute the ratio of makespans ALNS / competitor and report its mean with a
paired bootstrap 95\% confidence interval (10,000 resamples) \cite{efron1993introduction}. Each test is a one-sided
paired $t$-test of the mean log ratio below zero, with Holm's correction \cite{holm1979simple} within each of the three
families at $\alpha = 0.05$. C1 holds on a setting if all its tests are significant, and we claim C1 if it holds on at
least six of the eight settings. There is no minimum effect size; we report effect sizes throughout.

\paragraph{Pre-registration.} Hypotheses, competitors and their settings, budgets, instances, execution rules, analysis
and decision rule were written down and frozen before any ALNS or CP-SAT run on the test split. The frozen document
records the git object identifiers of the source trees, and our harness refuses to run on the test split unless the
checked-out code is exactly those trees. The test split was run once. Its known weaknesses are listed in it and in
Section~\ref{sec:discussion}. All seeds are fixed: one run per instance, method and budget.
